% Template from Overleaf https://www.overleaf.com/latex/templates/ieee-conference-template/grfzhhncsfqn 

\documentclass[conference]{IEEEtran}
\IEEEoverridecommandlockouts
% The preceding line is only needed to identify funding in the first footnote. If that is unneeded, please comment it out.
\usepackage{cite}
\usepackage{amsmath,amssymb,amsfonts}
\usepackage{algorithmic}
\usepackage{graphicx}
\usepackage{textcomp}
\usepackage{xcolor}
\usepackage{gensymb}
\pagestyle{plain}
\pagenumbering{arabic}
\def\BibTeX{{\rm B\kern-.05em{\sc i\kern-.025em b}\kern-.08em
    T\kern-.1667em\lower.7ex\hbox{E}\kern-.125emX}}
\begin{document}

\title{Horus: Dual-Polarimetric, Fully Digital, Software Defined Phased Array Radar -- The Perfect Successor for NEXRAD\\
%{\footnotesize \textsuperscript{*}Note: Sub-titles are not captured in Xplore and should not be used}
%\thanks{Identify applicable funding agency here. If none, delete this.}
}

\author{\IEEEauthorblockN{Christopher Lemus}
\IEEEauthorblockA{\textit{Department Of Electrical And Computer Engineering} \\
\textit{Boston University}\\
Boston, MA, United States Of America \\
Lemusc@bu.edu}

}

\maketitle

\begin{abstract}
On Earth we depend on reliable weather forecasts and reports on a day to day basis, and in some regions more than others becomes a matter of life or death. Some examples include communities living in tropical regions where large temperature gradients can evolve quickly and lead to severe showers, thunderstorms, and can kick off disasterous floods or mudslides\cite{b1}. Pilots on a day to day basis will note the forcast more than an hour ahead of their pre-flight checklist in order to coordinate the waypoints for the safest route for passenger and freight flight\cite{b1}. The key to providing the lastest information is to increase the frequency at which weather radars can provide imaging where necessary, hence the aim of this paper will be to provide facts supporting why we are ready to transition from the WSR-88D class weather surveillence radars in the United States of America (and even internationally) for a Dual-Polarimetric, Fully Digital, Software Defined Phased Array radar. 
\end{abstract}

\begin{IEEEkeywords}
radar, digital, par, nexrad, horus, simulation
\end{IEEEkeywords}
%%%%%%%%%%%%%%%%%%%%%%%%%%%%%%%%%%%%
%%%%%%%%%%%%%%%%%%%%%%%%%%%%%%%%%%%%
%%%%%%%%%%%%%%%%%%%%%%%%%%%%%%%%%%%%
%%%%%%%%%%%%%%%%%%%%%%%%%%%%%%%%%%%%
%%%%%%%%%%%%%%%%%%%%%%%%%%%%%%%%%%%%
\section{Introduction}
\subsection{History on Weather Radar}
During the Second World War a Scottish Physisict by the name of Sir Robert A. Watson-Watt applied his knowledge and experience in electromagnetics and sferics to detect aircraft at a distance\cite{history}. Naturally, discoveries made in such a fashion peak the curiosity of scientists and efforts to further enhance capability took place. By 1946 the organization at the time known as the Weather Bureau, now National Weather Service, rolled into effect the first set of radars in the Weather Surveillence Radar (WSR) family\cite{history}. A great majority of these radars operated in the X-band and few S-band were trickled in. These were some of the first to prove the concept of a network of weather radars providing data to the country. Fast forward a decade, the WSR-57 family of radars were rolled out with much greater promise versus the WSR-1, -2 and -3. These were also S-band radars with improved transmission power, gain, angular resolution and scanning time. A generation later, the modern day WSR-88D family radars were rolled out into the network of WSR's, which later earned the name of NEXRAD. This is the nework the National Weather Service (NWS), Federal Avaiation Adminsitration (FAA) and Department of Defense (DoD) have invested in to maintain the highest level measurement reliability on a day to day basis for civilian, aviation and military purposes to track weather across the United States\cite{history}. 
%%%%%%%%%%%%%%%%%%%%%%%%%%%%%%%%%%%%
%%%%%%%%%%%%%%%%%%%%%%%%%%%%%%%%%%%%
%%%%%%%%%%%%%%%%%%%%%%%%%%%%%%%%%%%%
%%%%%%%%%%%%%%%%%%%%%%%%%%%%%%%%%%%%
%%%%%%%%%%%%%%%%%%%%%%%%%%%%%%%%%%%%
\subsection{Goals for this paper}
The University of Oklahoma's (OU) Advanced Radar Research Center (ARRC) in conjunction with the National Oceanic and Atmospheric Administration (NOAA) have built one of the first dual-polarimetric, fully digital, software defined phased array radars for scientific research, called Horus\cite{b1}. There are many amazing features contained with this class of radar, including being fully digital, software defined, capable of multi-beam generation and reception, horizontal/vertical/circular polarization, transportability and scalability. In the paper by Robert D. Palmer et al., several flavors of analysis were provided spanning from calibration to real world ranging and detection performance of 20\% of the radar. 
\par The objective with this paper is to confirm their specifications for the proven 20\% of Horus, theoretical 100\% and theoretical 250\% of radar capability via MATLAB simulation. Parameters to be confirmed will be half power beam width (HPBW), element and array beam pattern, antenna gain and sensitivity. Presented will also be an extrapolation via scaling of antenna elements for the Horus radar to determine when it will match the specifications of the modern NEXRAD network radars. Ultimately, discovering and verifying what Horus brings to the table as the potential next generation system is the key driving factor in this report. The following sections will be laid out as follows. Section II will cover the theory \& methodology needed and used to confirm and verify the parameters shared between all radars. Section III will cover the results from simulation, useful figures to assist in visualization and drawing conclusions, along with a brief analysis of the data collected from the Palmer et al. journal article. Section IV will consist of summary of findings. Section V will conclude this  report and provide an answer to whether OU's Hours Phased Array Radar (PAR) is a proper candidate as successor for WSR-88D in the NEXRAD network. Acknowledgments, tables, figures and references will be at the end.
%%%%%%%%%%%%%%%%%%%%%%%%%%%%%%%%%%%%
%%%%%%%%%%%%%%%%%%%%%%%%%%%%%%%%%%%%
%%%%%%%%%%%%%%%%%%%%%%%%%%%%%%%%%%%%
%%%%%%%%%%%%%%%%%%%%%%%%%%%%%%%%%%%%
%%%%%%%%%%%%%%%%%%%%%%%%%%%%%%%%%%%%
\section{Theory \& Methodology}
\subsection{Element \& Array Beampattern Formulation}
\par The 2-Dimensional Antenna Pattern of any shaped aperture can be described as follows:
\begin{equation}
    E(\phi,\psi) = \frac{1}{\lambda^2}S_0(\frac{sin\phi cos\psi}{\lambda}, \frac{sin\phi sin\psi}{\lambda})
\end{equation}
where \emph{$S_0$} if the Fourier transform of the Illumination function, \emph{$s_0$}, which depends on the shape of the antenna element of interest\cite{blahut}. For the case of Horus PAR, circular antenna elements are used in an 8 x 8 lattice for each panel. There are 25 panels arranged in a 5 x 5 configuration. The illumination function is represented by
\begin{equation}
    s_0(x,y) = circ(\frac{x}{2r}, \frac{y}{2r})
\end{equation}
who has the following Fourier transform:
\begin{equation}
    S_0(x,y) = 2r\cdot2rjinc(\frac{2r}{\lambda}sin\phi) = 4r^2jinc(\frac{2r}{\lambda}sin\phi) 
\end{equation}
therefore the antenna pattern for a single element is:
\begin{equation}
    E(\phi,\psi) = \frac{4r^2}{\lambda^2}jinc(\frac{2r}{\lambda}sin\phi) 
\end{equation}
In order to represent any N x M array pattern we have to account for each row and column. The principle array pattern formula is:
\begin{equation}
    E(\phi,\psi) = E_e(\phi,\psi)dirc_N(\frac{d_1}{\lambda}sin\phi cos\psi)dirc_M(\frac{d_2}{\lambda}sin\phi sin\psi)
\end{equation}
where \emph{$E_e$} is equal to equation (4), the element spacing \emph{$d_1$} = \emph{$d_2$} = $\frac{\lambda}{2}$ meters. To further simplify (5) we can futher expand the Dirichlet functions \emph{$dirc_N$} and \emph{$dirc_M$} as follows:
 \begin{equation}
    dirc_N(\frac{d_1}{\lambda}sin\phi cos\psi) = \frac{sin[N\pi\frac{d_1}{\lambda}sin\phi cos\psi]}{sin[\pi\frac{d_1}{\lambda}sin\phi cos\psi]}
\end{equation}
 \begin{equation}
    dirc_M(\frac{d_2}{\lambda}sin\phi sin\psi) = \frac{sin[M\pi\frac{d_2}{\lambda}sin\phi sin\psi]}{sin[\pi\frac{d_2}{\lambda}sin\phi sin\psi]}
\end{equation}
Therefore the holy grail equation to describe the antenna pattern is 
\begin{equation}
    E(\phi,\psi) = E_e(\phi,\psi)\frac{sin[N\pi\frac{d_1}{\lambda}sin\phi cos\psi]}{sin[\pi\frac{d_1}{\lambda}sin\phi cos\psi]}\frac{sin[M\pi\frac{d_2}{\lambda}sin\phi sin\psi]}{sin[\pi\frac{d_2}{\lambda}sin\phi sin\psi]}
\end{equation}
\par The power pattern is defined as the square of the antenna element pattern multiplied by the number of elements with respect to the steering angle $\phi$:
\begin{equation}
    F(\phi) = |E_e(\phi)|^2dirc_N^2(\frac{d}{\lambda}sin\phi - sin\phi_0)
\end{equation}
Equation (9) will also be used to show what the HPBW is for a given N x M PAR throughtout our evaluations. 
%%%%%%%%%%%%%%%%%%%%%%%%%%%%%%%%%%%%
%%%%%%%%%%%%%%%%%%%%%%%%%%%%%%%%%%%%
%%%%%%%%%%%%%%%%%%%%%%%%%%%%%%%%%%%%
%%%%%%%%%%%%%%%%%%%%%%%%%%%%%%%%%%%%
%%%%%%%%%%%%%%%%%%%%%%%%%%%%%%%%%%%%
\subsection{Antenna Gain}
\par Much of radar design is centered around the antenna chosen because that is the interface between the world and electronics we model. Naturally it is important to care about antenna gain, which according to sheet 106 from Kudeki's lecture notes \cite{kudeki} can be represented as follows:
\begin{equation}
    G = (\frac{4\cdot\pi\cdot A_e}{\lambda^2})
\end{equation}
where \emph{$A_e$} equals the antenna effective area. More formally the gain is represented in decibels as the values can get quite large:
\begin{equation}
    G_{dB} = 10\cdot \text{log}_{10}(\frac{4\cdot\pi\cdot A_e}{\lambda^2})
\end{equation}

%%%%%%%%%%%%%%%%%%%%%%%%%%%%%%%%%%%%
%%%%%%%%%%%%%%%%%%%%%%%%%%%%%%%%%%%%
%%%%%%%%%%%%%%%%%%%%%%%%%%%%%%%%%%%%
%%%%%%%%%%%%%%%%%%%%%%%%%%%%%%%%%%%%
%%%%%%%%%%%%%%%%%%%%%%%%%%%%%%%%%%%%
%%%%%%%%%%%%%%%%%%%%%%%%%%%%%%%%%%%%
%%%%%%%%%%%%%%%%%%%%%%%%%%%%%%%%%%%%
%%%%%%%%%%%%%%%%%%%%%%%%%%%%%%%%%%%%
%%%%%%%%%%%%%%%%%%%%%%%%%%%%%%%%%%%%
%%%%%%%%%%%%%%%%%%%%%%%%%%%%%%%%%%%%
\subsection{Radar Resolution}
\par The range resolution for a linearly modulated pulse compressed electromagnetic wave can be represented by:
\begin{equation}
    R_r = \frac{c}{2B} = \frac{c\cdot\tau}{2B}
\end{equation}
where \emph{c} is the speed of light, \emph{$\tau$} is pulse width and \emph{B} is the receiver bandwidth. In the case of the Horus PAR the max pulse width is 100$\mu$s at 10\% duty cycle and therefore can match the range resulution of the WSR-88D by using the same pulse width. If we calculate the range resolution with the WSR-88D short pulse specification we see it has a range resolution of 250 meters.
\subsection{Doppler \& Polarization Base Products}
\par For a radar with no capability to alternate between different types of polarization, generally the data is solely Dopper based\cite{wolffDualPol}. If there is polarimetric capability (in this case dual) then polarization based products can be derrived. In particaular we would like to briefly discuss the types of Doppler and polarization products that NEXRAD radars and Horus generate. A comparison of the data from Palmer et al. \cite{b1} will be provided in the Results section. 
\par The Doppler based products are Radar Reflectivity ($Z$ in dBZ), Doppler Velocity ($v$ in ms$^{-1}$) and Spectrum Width ($
\sigma_v$ in ms$^{-1}$). The polarization based products are Differential Reflectivity (Z$_{DR}$ in dB), Differential Phase ($\phi_{DR}$ in degrees) and Correlation Coefficient ($\rho_{hv}$ in dimensionless units)
\par The Radar Reflectivity is a gauge for how strong a reflector is in a volume. The Doppler Velocity is an indicator of speed and directionality of objects in a volume. The Spectrum Width is the vaiance of velocities within the volume that is scanned, which tells us a lot of information about what is going on in a storm. Differential polarization is defined as the difference between the reflectivites of Horizonal and Vertical electromagnetic wave reflections. This is valuable because the structure and properties of the interrogated volume and hydrometeors will cause a different response from the electromagnetic wave. Different sensed magnitudes of reflectivites are associated with different polatizations. Similarly, the Differential Phase is the difference between the phases of Horizontal and Vertical electromagnetic pulses. These values will also vary drastically depending on the structure and properties of the interrogated volume and hydrometeors. Finally, the Correlation Coefficient defines how uniform particles are within the interrogated volume. These are all base prodcts and can be used to derive what are known as "derrived products" which are used to describe other properties of a storm.
%%%%%%%%%%%%%%%%%%%%%%%%%%%%%%%%%%%%
%%%%%%%%%%%%%%%%%%%%%%%%%%%%%%%%%%%%
%%%%%%%%%%%%%%%%%%%%%%%%%%%%%%%%%%%%
%%%%%%%%%%%%%%%%%%%%%%%%%%%%%%%%%%%%
%%%%%%%%%%%%%%%%%%%%%%%%%%%%%%%%%%%%
\subsection{Methodology}
\par MATLAB is used for simulation to prove the results and claims that were made for the Horus PAR. Scripts were created for each of the 3 cases (320, 1600, 10000 elements) to show the unique array beam patterns (assuming no multi-beam) and HPBW. Implementing the expressions described in the previous subsections above we are able to generate Figures 1-9 which are provided and analyzed in the Results section. Parameters such as number of elements in array, seperation distance between elements, individual antenna element beampattern, operating frequency/wavelength can be updated in the script to analyze other systems, but for this paper only number of elements were adjusted and all other parameters remain constant. We were also able to use the known Horus parameters to calculate the Antenna area in square meters and Antenna Gain. 
%%%%%%%%%%%%%%%%%%%%%%%%%%%%%%%%%%%%
%%%%%%%%%%%%%%%%%%%%%%%%%%%%%%%%%%%%
%%%%%%%%%%%%%%%%%%%%%%%%%%%%%%%%%%%%
%%%%%%%%%%%%%%%%%%%%%%%%%%%%%%%%%%%%
%%%%%%%%%%%%%%%%%%%%%%%%%%%%%%%%%%%%
\section{Results}
\par  In the journal by Palmer et al. \cite{b1} claims of beamwidths of 13$^\circ$ in azimuth by 3.1$^\circ$ in elevation are made for the 5x1 array configuration they tested. They specify the fully populated radar can achieve a 2.58$^\circ$ beamwidth in azimuth and elevation, and at 10,000 elements a beamwidth $\approx$1$^\circ$. Figures 1 shows the configuration under which results were gathered by the team at OU's ARRC. Figure 4 highlights the beampattern for that configuration and Figure 7 shows the proof that the beamwiths specified are attainable. Figure 2 shows the configuration of 5x5 which the team did not get to collect data for because of lack of hardware to populate the other 4 colums of the PAR. Figure 5 shows the beampatter for that configuration, followed by Figure 8 which confirms the half power beamwidth claimed in their specification. Lastly, the more theoretical 10,000 element case is shown as well in Figures 3, 6 and 9, which highlight the feasability of a 1 degree beamwidth. The significance of these findings are that at scale a phased array of this caliber can be a direct replacement for the WSR-88D.
\par Equally as important is the size of the PAR. Horus is specified to be a 2.03 m x 2.03 m array, totalling 4.121 square meters in area. The theoretical 10,000 element case, assuming a square 100 x 100 element configuration, would total to about 25.756 square meters. The WSR-88D has a radius of 4.2 meters, and being a dish has an area of 55.43 square meters. This means the PAR would be physically a smaller radar with equal, if not more, capability. The current NEXRAD radars could theoretically be swapped directly with a 10,000 element PAR with room to spare in the radome. Of course, in practice, roadblocks are guaranteed. 
\par Another key benefit is that the fully digital, software defined aspect allows for significantly quicker scan times. In context of the Palmer et al. journal, Range Height Indicator (RHI) is used, which means a fixed Azimuth angle ange varying Elevation angles from 0.5$^\circ$ to 32.5$^\circ$ in 0.5$^\circ$ incrememnts. At each Elevation, there are 64 samples collected, each taking 1 ms of time to complete. This means for all elevations at one Azimuth it takes 4.096 Seconds to scan and can conduct a full Vertical Coverage Pattern (VCP) scan in 10-12 seconds\cite{b4} using muli-beam generation\cite{b6}. The WSR-88D for comparison has a VCP ranging from 0.5$^\circ$ to 19.5$^\circ$ and takes anywhere from 4 to 6 minutes\cite{b5} using only a single beam.
\par The final key point to highlight is the ability for the PAR to collect polarization product information due to its dual polarimetric capability. We present the results from the journal in Figure 10, with Doppler products above and Polarimetric products below. The Horus radar operating at 20\% of its full capability captured the results in its RHI mode. It was not conductiong measurements at 360$^\circ$ but at a fixed azimuth. The NEXRAD WSR-88D at KTLX was operating in the PPI scanning mode which is why they were able to collect data from 0 to 360$^\circ$ azimuth. We can see that the results line up very nice with the NEXRAD data captured at the same time frame as the Hours PAR. 
\par Table I sumarizes the key parameters for the latest WSR-88D, 20\% Horus capability corresponding to 320 antenna elements in 40x8 array configuration, 100\%  Hours capability corresponding to all 1,600 elements in the 40x40 array configuration and 250\% Horus capability correponding to 10,000 antenna elements in a 100x100 array configuration. The parameters in this table were consolodated from the Horus publication\cite{b1}, WSR-88D technical specifications\cite{b4} and deduced from the given parameters for the theoretical cases. 

%%%%%%%%%%%%%%%%%%%%%%%%%%%%%%%%%%%%
%%%%%%%%%%%%%%%%%%%%%%%%%%%%%%%%%%%%
%%%%%%%%%%%%%%%%%%%%%%%%%%%%%%%%%%%%
%%%%%%%%%%%%%%%%%%%%%%%%%%%%%%%%%%%%
%%%%%%%%%%%%%%%%%%%%%%%%%%%%%%%%%%%%
\section{Summary}
\par The work done by ARRC at OU is pivital because the team was able to correlate and validate their results in a very direct way with the pre-existing sensors in the NEXRAD network. This serves as a great benchmark and is certain to have provided insight for improvement during runtime. They are on the correct trajectory to constructing a fully operational Horus PAR which can provide the public with fast, reliable data and the research community with flexability to experiemnt with nonlinear frequency modulation schemes to improve resolution.
%%%%%%%%%%%%%%%%%%%%%%%%%%%%%%%%%%%%
%%%%%%%%%%%%%%%%%%%%%%%%%%%%%%%%%%%%
%%%%%%%%%%%%%%%%%%%%%%%%%%%%%%%%%%%%
%%%%%%%%%%%%%%%%%%%%%%%%%%%%%%%%%%%%
%%%%%%%%%%%%%%%%%%%%%%%%%%%%%%%%%%%%
\section{Conclusion}
\par In it's current state, the Horus cannot directly replace the WSR-88D class radar. However, if the number of elements reach approximately 10,000 then the specifications become much more comparable and at that point can certainly become the successor. It's clear that the digital nature of the phased array is a major benefit and tons of data becomes readily available at a significantly faster rate. The significance of the dual polarimetric capability was shown to yield promising results by direct comparison to the NEXRAD dataset.
%%%%%%%%%%%%%%%%%%%%%%%%%%%%%%%%%%%%
%%%%%%%%%%%%%%%%%%%%%%%%%%%%%%%%%%%%
%%%%%%%%%%%%%%%%%%%%%%%%%%%%%%%%%%%%
%%%%%%%%%%%%%%%%%%%%%%%%%%%%%%%%%%%%
%%%%%%%%%%%%%%%%%%%%%%%%%%%%%%%%%%%%
\section*{Acknowledgment}

I would like to thank Dr. Joshua Semeter for offering an excellent EC707 Radar Remote Sensing course at Boston University and being thoughroughly invested in his students. I would also like to thank Boston University for allowing me the opportunity and wherewithal to be a contributor to larger-than-life missions. Equally as important, I would like to thank The University of Oklahoma's Adanced Radar Research Center (ARRC) and all co-contributors to this excellent, masterful feat of engineering, and offering an student/engineer like myself the resources necessary to analyze their complex system and it's performace. Their work is truly paving the way towards a better future.
%%%%%%%%%%%%%%%%%%%%%%%%%%%%%%%%%%%%
%%%%%%%%%%%%%%%%%%%%%%%%%%%%%%%%%%%%
%%%%%%%%%%%%%%%%%%%%%%%%%%%%%%%%%%%%
%%%%%%%%%%%%%%%%%%%%%%%%%%%%%%%%%%%%
%%%%%%%%%%%%%%%%%%%%%%%%%%%%%%%%%%%%
\begin{thebibliography}{00}
\bibitem{b1} Palmer, Robert D, et al. Horus—a Fully Digital Polarimetric Phased Array Radar for Next-Generation Weather Observations. Vol. 1, 1 Jan. 2023, pp. 96–117, https://doi.org/10.1109/trs.2023.3280033. Accessed 26 Sept. 2023.
\bibitem{history} Whiton, Roger C., et al. “History of Operational Use of Weather Radar by U.S. Weather Services. Part I: The Pre-NEXRAD Era.” Weather and Forecasting, vol. 13, no. 2, 1 June 1998, pp. 219–243
\bibitem{b3} “Defense Technical Information Center.” Dtic.mil, 2026, apps.dtic.mil/sti/citations/ADA273077. Accessed 15 Apr. 2026.
\bibitem{b5}NEXRAD Technical Information. https://www.roc.noaa.gov/public-documents/wsr88d/NEXRAD-Technical-Information.pdf
\bibitem{b4} Schultz, Cory, et al. Investigation of Lightning and Storm Electrification Processes Using a Phased Array Radar and Lightning Mapping Array.
\bibitem{b6} Advanced Weather Surveillance Capabilities of the Fully Digital Horus Phased Array Radar, journals.ametsoc.org/view/journals/bams/107/1/BAMS-D-24-0296.1.xml. Accessed 15 Apr. 2026. 
\bibitem{b7} Mahafza, Bassem R. Radar Systems Analysis and Design Using MATLAB. CRC Press, 2000. 
\bibitem{blahut} “Antenna Systems.” Theory of Remote Image Formation, 18 Nov. 2004, pp. 153–185, https://doi.org/10.1017/cbo9780511543418.006.
\bibitem{kudeki} ECE 458 Lecture Notes on Applications of Radiowave Propagation Erhan Kudeki, courses.physics.illinois.edu/ece458/458sp10.pdf. Accessed 1 May 2026.  
\bibitem{NOAAsite} “NOAA’s Weather and Climate Toolkit.” National Centers for Environmental Information (NCEI), 21 Nov. 2025, www.ncei.noaa.gov/products/weather-climate-toolkit. 
\bibitem{wolffDualPol} Wolff, Dipl.-Ing. (FH) Christian. “Radar Basics.” Radartutorial, www.radartutorial.eu/15.weather/wr51.en.html. Accessed 1 May 2026. 
\end{thebibliography}
%%%%%%%%%%%%%%%%%%%%%%%%%%%%%%%%%%%%
%%%%%%%%%%%%%%%%%%%%%%%%%%%%%%%%%%%%
%%%%%%%%%%%%%%%%%%%%%%%%%%%%%%%%%%%%
%%%%%%%%%%%%%%%%%%%%%%%%%%%%%%%%%%%%
%%%%%%%%%%%%%%%%%%%%%%%%%%%%%%%%%%%%

\section*{Figures \& Tables}
\begin{figure}[htbp]
    \centering
    \includegraphics[width=0.3\textwidth]{figs/array5x1.png}
    \caption{Topological view of 40 x 8 radiating elements}
    \label{view320}
\end{figure}
\begin{figure}[htbp]
    \centering
    \includegraphics[width=0.3\textwidth]{figs/array5x5.png}
    \caption{Topological view of 40 x 40 radiating elements}
    \label{view1600}
\end{figure}
\begin{figure}[htbp]
    \centering
    \includegraphics[width=0.3\textwidth]{figs/array10000.png}
    \caption{Topological view of 100 x 100 radiating elements}
    \label{view100000}
\end{figure}
\begin{figure}[htbp]
    \centering
    \includegraphics[width=0.5\textwidth]{figs/bp320.png}
    \caption{Array pattern of 40 x 8 radiating elements}
    \label{arrayPattern320}
\end{figure}
\begin{figure}[htbp]
    \centering
    \includegraphics[width=0.5\textwidth]{figs/bp1600.png}
    \caption{Array pattern of 40 x 40 radiating elements}
    \label{arrayPattern1600}
\end{figure}
\begin{figure}[htbp]
    \centering
    \includegraphics[width=0.5\textwidth]{figs/bp10000.png}
    \caption{Array pattern of 100 x 100 radiating elements}
    \label{arrayPattern10000}
\end{figure}\begin{figure}[htbp]
    \centering
    \includegraphics[width=0.5\textwidth]{figs/hpbw5x1.png}
    \caption{Half power beamwidth of 40 x 8 radiating elements}
    \label{arrayPattern320}
\end{figure}
\begin{figure}[htbp]
    \centering
    \includegraphics[width=0.5\textwidth]{figs/hpbw5x5.png}
    \caption{Half power beamwidth of 40 x 40 radiating elements}
    \label{arrayPattern1600}
\end{figure}
\begin{figure}[htbp]
    \centering
    \includegraphics[width=0.5\textwidth]{figs/hpbw10000.png}
    \caption{Half power beamwidth of 100 x 100 radiating elements}
    \label{arrayPattern10000}
    \centering
    \includegraphics[width=0.5\textwidth]{figs/journalData.png}
    \caption{Results of data captured from KTLX and Horus on 23 March 2023 20:35:06 Z}
    \label{arrayPattern10000}
\end{figure}

\vspace{12pt}
\begin{table*}[htbp]
\caption{Benchmark Parameters for Horus to Meet}
%\begin{center}
\begin{tabular}{|c|c|c|c|c|}
\hline
\cline{2-4} 
\textbf{Parameter} & \textbf{\textit{WSR-88D (Proven)}}& \textbf{\textit{20\% Horus (Proven)}}& \textbf{\textit{100\% Hours (Theory)}} & \textbf{\textit{250\% Horus (Theory)}} \\
\hline
Operating Frequency ($GHz$)&2.7-3.0&\multicolumn{3}{c|}{2.7-3.1}\\
\hline
Element Polarization&Dual Polarization&\multicolumn{3}{c|}{ATSR/STSR/LHCP/RHCP}\\
\hline
Tx Waveform Type&P0N&\multicolumn{3}{c|}{AWG/LFM/NLFM}\\
\hline
Tx Pulse Width Short ($\mu$s)&1.52-1.65&\multicolumn{3}{c|}{100 max}\\
\hline
Tx Pulse Width Long ($\mu$s)&4.61-4.81&\multicolumn{3}{c|}{100 max}\\
\hline
Tx Beamwidth ($degrees$)&0.925&13 x 2.58&2.58&0.99\\
\hline
Tx Bandwidth ($MHz$)&72.6&\multicolumn{3}{c|}{100}\\
\hline
Tx Gain ($dB$)&48 dB&30.15&37.15&45.1\\
\hline
Antenna Area ($m^2$)&55.43&0.824&4.121&25.756\\
\hline
Beamsteering Range (az/el $degrees$)&-3 to 65&\multicolumn{3}{c|}{-45 to 45}\\
\hline
Mechanical position (az/el $degrees$)&360 / -3 to 65&\multicolumn{3}{c|}{360 / -1 to 92}\\
\hline
Scanning time (minutes)&4-6&\multicolumn{3}{c|}{1 to 1.5}\\
\hline
\multicolumn{5}{l}{$^{\mathrm{a}}$Horus Data from \cite{b1} and WSR-88D data from \cite{b5}}
\end{tabular}
\label{tab1}
%\end{center}
\end{table*}
%%%%%%%%%%%%%%%%%%%%%%%%%%%%%%%%%%%%
%%%%%%%%%%%%%%%%%%%%%%%%%%%%%%%%%%%%
%%%%%%%%%%%%%%%%%%%%%%%%%%%%%%%%%%%%
%%%%%%%%%%%%%%%%%%%%%%%%%%%%%%%%%%%%
%%%%%%%%%%%%%%%%%%%%%%%%%%%%%%%%%%%%




\end{document}
